\section{Introduction}
\label{sec:intro}

The deployment of mixed-criticality workloads on shared hardware platforms is
a central challenge in modern embedded and cyber-physical systems.
Hypervisors provide spatial and temporal partitioning between guest domains of
different criticality levels, making them an attractive building block for
consolidating safety-critical and best-effort applications on a single physical
machine.
However, the scheduling overhead, interrupt handling, and resource contention
inherent in general-purpose hypervisors can compromise the determinism required
by hard real-time workloads.

Two of the most widely studied general-purpose hypervisors, \textbf{Xen} and
\textbf{KVM}, differ fundamentally in their architecture.
Xen is a Type-1 (bare-metal) hypervisor with a dedicated privileged domain
(Dom0) that mediates all device access; KVM is a Type-2 loadable kernel module
that extends a standard Linux host into a hypervisor.
Understanding how these architectural differences manifest as latency and jitter
under realistic workloads is essential for practitioners who must select and
configure a virtualization stack for mixed-criticality deployments.


\subsection{Related Work}
\label{sec:related}

Abeni and Faggioli~\cite{abeni2019} conducted a systematic experimental
analysis of Xen and KVM latency profiles, characterizing the impact of the
\texttt{PREEMPT\_RT} patch at both host and guest level and quantifying the
scheduling interference introduced by each hypervisor under background stress
workloads.
Their follow-up work~\cite{abeni2020} extended this analysis to include
isolation mechanisms such as vCPU pinning and explored the Xen Null scheduler
as a means of achieving near-static partition allocation with near-zero
scheduling overhead.
A further study by the same author~\cite{abeni2024} examined the behavior of
virtualized real-time workloads in container-based and VM-based environments,
broadening the comparison to include lightweight virtualization alternatives.

Cinque~\emph{et al.}~\cite{cinque2024} assessed temporal isolation in
safety-critical mixed-criticality systems running on a Xen hypervisor,
identifying residual interference paths that can violate isolation boundaries
even after applying standard hardening techniques such as IRQ affinity and
VCPU pinning.
Their results highlight the importance of measuring not just average-case
but worst-case latency behavior across long observation windows.

% On the unikernel front, Kuenzer~\emph{et al.}~\cite{kuenzer2021} introduced
% Unikraft as a high-performance, specialized unikernel framework designed for
% minimal boot time and resource footprint.
% Unikraft allows applications to be compiled as standalone bootable images
% without a full operating system, making it a promising guest alternative
% for latency-sensitive workloads.


\subsection{Contributions and Experimental Scope}
\label{sec:contributions}

This work reproduces and extends the experimental campaign of
Abeni and Faggioli~\cite{abeni2019,abeni2020} on a current software stack
(Linux~6.x, Xen~4.17, mainline \texttt{PREEMPT\_RT} rt5).
Our specific contributions are:

\begin{itemize}
  \item \textbf{Baseline reproduction:} We replicate the key scheduling-latency
        measurements of~\cite{abeni2019} using \texttt{cyclictest} across
        all four host--guest kernel combinations
        (NRT$\times$NRT, NRT$\times$RT, RT$\times$NRT, RT$\times$RT),
        both with and without a background \texttt{stress-ng} workload,
        for both KVM and Xen.

  \item \textbf{Extended isolation analysis:} We evaluate the incremental
        effect of vCPU pinning, host-side \texttt{isolcpus}/\texttt{nohz\_full},
        IRQ affinity steering, and cache/memory-bandwidth partitioning
        via Intel Resource Director Technology (RDT).

  \item \textbf{Xen guest-type and scheduler comparison:} We compare all three
        Xen guest types---HVM, PV, and PVH---under the Credit2 and Null
        schedulers to identify the configuration that minimises worst-case
        latency.

  \item \textbf{TACLe benchmark experiments:} We execute five real-time
        workloads from the TACLeBench suite~\cite{taclebench}
        (DEBIE, \texttt{huff\_enc}, \texttt{lift}, \texttt{matrix1},
        \texttt{test3}) inside a critical guest domain alongside a noisy
        neighbour, quantifying worst-case execution time (WCET) degradation
        under each configuration.

%   \item \textbf{Unikraft unikernel evaluation:} We boot a
%         Unikraft~\cite{kuenzer2021} unikernel as a guest on both KVM and Xen
%         and compare its latency profile and resource footprint against a full
%         Linux \texttt{PREEMPT\_RT} guest.
\end{itemize}


\subsection{Experimental Platform}
\label{sec:platform}

All experiments were conducted on the same physical host equipped with Intel
VT-x hardware virtualization support and Intel Resource Director Technology
(RDT) for cache and memory-bandwidth partitioning.
The software stack comprised Xen~4.17-amd64 and KVM on a Linux~6.x host.
Guest kernels spanned three configurations: non-real-time (NRT, vanilla
Linux), Low-Latency (LL), and fully preemptible (\texttt{PREEMPT\_RT} rt5).
Scheduling latency was measured using \texttt{cyclictest}, and worst-case
execution time was characterized using the TACLeBench~\cite{taclebench} suite.
Background interference was generated using \texttt{stress-ng} with
cache-pressure, interrupt-flooding, and raw-socket stressors.

The remainder of this paper is organized as follows.
Section~\ref{sec:setup} describes the testbed configuration and kernel build
procedure.
Sections~\ref{sec:kvm} and~\ref{sec:xen} present the KVM and Xen experimental
results, respectively.
Section~\ref{sec:conclusion} concludes the paper and outlines future work.
